\section{三种反馈方法的系统比较}

\begin{table}[H]
\centering
\small
\begin{tabular}{>{\raggedright\arraybackslash}p{2.3cm}>{\raggedright\arraybackslash}p{3.4cm}>{\raggedright\arraybackslash}p{3.2cm}>{\raggedright\arraybackslash}p{4.2cm}}
\toprule
\textbf{方法} & \textbf{反馈来源} & \textbf{主要用途} & \textbf{关键风险} \\
\midrule
ReAct & 搜索、网页、工具 observation & 当前轨迹规划与纠错 & 不可信内容、循环、状态漂移 \\
RLEF & 编译器、测试、运行日志 & 多轮修复与 RL 更新 & 测试投机、沙箱 exploit、覆盖不足 \\
Constitutional AI & 原则驱动的 AI critique / preference & SFT 与 RLAIF & 原则冲突、judge 偏差、过度拒答 \\
\bottomrule
\end{tabular}
\end{table}

\subsection{共同架构}

三者都可以拆成：

$$
\text{proposal}\rightarrow\text{feedback function}
\rightarrow\text{revised proposal}\rightarrow\text{selection / update}.
$$

区别在于 feedback function 的客观性和成本。执行反馈最接近 ground truth，但只适用于可执行任务；AI feedback 覆盖面广，却需要更强外部校准。

\subsection{设计原则}

\begin{enumerate}
    \item 能用真实环境验证时，优先环境信号；
    \item 将反馈分解到可行动的错误粒度；
    \item 把不可逆动作放在权限边界之外；
    \item 训练 reward 与最终评估指标保持独立；
    \item 记录所有反馈来源与模型版本，支持回放。
\end{enumerate}

\subsection{本章小结}

反馈越自动化，系统越容易扩展；反馈越接近真实结果，系统越不易漂移。高质量自我改进系统通常混合多种信号：外部执行作为锚点，模型 critique 提供密集解释，人类负责原则与高风险边界。

